\section{Fonctionnement de \texorpdfstring{\taub}{tau2-bench}}
\label{sec:benchmark}

\subsection{Une simulation à trois acteurs}

Chaque épisode met en jeu trois composants, coordonnés par un orchestrateur qui régit les tours d'interaction :

\vspace{2mm}
\begin{description}[nosep]
  \item[L'agent évalué] est le modèle testé. Son prompt système intègre la politique métier du domaine en langue naturelle (définissant les règles de calcul et les conditions d'éligibilité et restrictions) ainsi que les spécifications des outils mis à sa disposition.
  \item[L'utilisateur simulé] est un second LLM paramétré par un profil privé comprenant persona, objectif précis, éléments d'information connus ou ignorés et consignes comportementales (\emph{e.g.} fermeté, négociation, réticence).
  \item[L'environnement] est composé d'une base de données relationnelle en mémoire et d'outils Python déterministes modélisant les actions de consultation et de modification de l'état système.
\end{description}

\vspace{2mm}\noindent
L'épisode s'achève lors de l'émission d'un token d'arrêt explicite par l'utilisateur (\texttt{\#\#\#STOP\#\#\#}), basculement vers un opérateur humain, ou encore un dépassement du plafond de tours autorisé (\texttt{max\_steps}).

\subsection{Cadre décisionnel et extension en contrôle partagé (\emph{dual-control})}

Alors que les domaines classiques ($\tau$-bench) modélisent l'interaction comme un processus de décision markovien partiellement observable (POMDP) où seul l'agent dispose d'outils sur l'environnement, le domaine \texttt{telecom} de \taub{} introduit une dynamique décentralisée formalisée sous l'angle d'un Dec-POMDP \cite{barres2025tau2}. 

\vspace{2mm}\noindent
Dans ce cadre, l'utilisateur possède son propre état privé (son téléphone mobile : statut réseau, mode avion, configurations APN) et une palette d'outils dédiés (30 outils utilisateurs). L'agent de support n'a aucun accès direct à cet état terminal et ne peut le modifier : il doit guider l'utilisateur par le dialogue afin que ce dernier exécute les bonnes manipulations techniques et lui rapporte ses observations. Pour l'agent, cette configuration transforme le problème en un suivi de croyance (\emph{belief tracking}) sous incertitude asymétrique : l'agent doit raisonner sur ce que l'utilisateur voit et comprend afin d'ordonnancer correctement ses instructions.

\subsection{Métrique de fiabilité : \texorpdfstring{$\passk{k}$}{pass\^{}k}}

Chaque essai d'une tâche reçoit une récompense binaire $r \in \{0,1\}$ (section~\ref{sec:evaluation}). Sur $n$ essais indépendants d'une tâche générant $c$ réussites, l'estimateur sans biais de la probabilité que $k$ essais consécutifs réussissent tous est :
\begin{equation}
  \passk{k} = \mathbb{E}_{\text{tâches}}\left[\frac{\binom{c}{k}}{\binom{n}{k}}\right], \qquad 1 \le k \le n.
\end{equation}
Tandis que $\passk{1}$ correspond au taux de réussite moyen classique, $\passk{k}$ décroît strictement avec $k$. Pour un agent déployé en production, $\passk{2}$ et $\passk{3}$ quantifient la \emph{reproductibilité} des décisions face à des variations stochastiques du dialogue.

\subsection{Modes d'exécution et isolation des facteurs}

Le benchmark propose plusieurs modes d'exécution qui permettent d'isoler expérimentalement les sources d'échec :
\begin{itemize}[nosep]
  \item \textbf{Mode standard (conversationnel)} : interaction bidirectionnelle complète entre l'agent et l'utilisateur simulé.
  \item \textbf{Mode solo (sans utilisateur)} : l'agent reçoit l'intégralité du scénario utilisateur dans son prompt et doit générer la séquence d'outils appropriée sans interaction interactive. Ce mode permet de découpler la capacité intrinsèque de raisonnement et de planification de l'agent des perturbations conversationnelles
  \item \textbf{Mode plan oracle} : l'agent est guidé par le plan d'actions idéal pour mesurer sa seule fidélité d'exécution syntaxique.
\end{itemize}
